\documentclass[11pt,a4paper]{article}
\usepackage[T1]{fontenc}
\usepackage[utf8]{inputenc}
\usepackage[english]{babel}
\usepackage{lmodern,amsmath,amssymb,amsthm,mathtools,mathrsfs}
\usepackage[margin=1.02in]{geometry}
\usepackage{microtype,enumitem,xurl}
\usepackage[hidelinks]{hyperref}
\numberwithin{equation}{section}
\newtheorem{theorem}{Theorem}[section]
\newtheorem{lemma}[theorem]{Lemma}
\newtheorem{proposition}[theorem]{Proposition}
\theoremstyle{remark}\newtheorem{remark}[theorem]{Remark}
\newcommand{\N}{\mathbb N}
\newcommand{\Z}{\mathbb Z}
\newcommand{\R}{\mathbb R}
\newcommand{\T}{\mathbb R/\mathbb Z}
\newcommand{\1}{\mathbf 1}
\newcommand{\ee}{\mathrm e}
\newcommand{\E}{\mathbb E}
\DeclareMathOperator{\Tr}{Tr}
\emergencystretch=2em
\setlist{itemsep=3pt,topsep=5pt}
\title{From the Liouville resonance remainder\\to ordinary two-point cancellation}
\author{A remainder-first derivation from V. Casa\~na-Giner's decomposition\\
using intermediate estimates of the supplied OpenAI manuscript}
\date{7 November 2026}
\begin{document}
\maketitle
\begin{abstract}
The starting point is the nonproportional block remainder in
\emph{Rational Resonances and Microscopic Natural--Logarithmic Profiles for
the Liouville Function}. We define those blocks before taking their sum,
prove their exact conversion to progression intervals, and close the
remainder by a centered-divisor and graph argument. The final affine
correlation theorem of the supplied \texttt{final.pdf} is never used as
an input. Instead, specific intermediate Fourier, deletion, and graph-moment
estimates from that manuscript are stated as external inputs. Their
consequences are developed here, including the operator transfer and
the finite-interval boundary estimates. Finally, we recover ordinary
two-point Chowla directly from the negative microscopic remainders,
by an explicit induction on the shift.

This is a proof relative to the intermediate results explicitly cited
below. It is not a self-contained reproof of their sieve, circuit, or
closed-walk arguments, and it does not establish an independent proof
of those results. No source sections are appended or included.
\end{abstract}

\paragraph{Provenance.}
\textbf{[LIT]} identifies results imported from the supplied sources.
\textbf{[ADAPT]} identifies their application to the resonance remainder.
\textbf{[DED]} identifies derivations written out here. These labels describe
logical provenance, not a claim of historical priority. In particular,
the negative-remainder extraction in Section~\ref{sec:chowla} is a
deduction here; no literature-wide novelty claim is made for it.

\tableofcontents

\section{Start with the nonproportional remainder}
\label{sec:remainder}

Write $\lambda(n)=(-1)^{\Omega(n)}$ for $n\ge1$, so that
$\lambda(mn)=\lambda(m)\lambda(n)$ and $\lambda(n)^2=1$.
For fixed coprime $p,q\in\N$ and fixed $c\in\R$, put
\begin{equation}\label{eq:CX}
 C_X(u)=\sum_{\substack{1\le n\le X\\\lfloor e^u n\rfloor\ge1}}
 \lambda(n)\lambda(\lfloor e^u n\rfloor),\qquad
 u_X=\log(p/q)+c/X.
\end{equation}
We initially take $X\in\N$. Define
\[
 b_a=\lfloor pa/q\rfloor,\qquad
 s_a=pa-qb_a\in\{0,\ldots,q-1\},\qquad
 \delta_X=\frac pq(e^{c/X}-1).
\]
For $0\le a<q$ and $r\in\Z$, let $I_{a,r}(X,c)$ be the set of integers
$k\ge0$ satisfying all three restrictions
\begin{equation}\label{eq:I}
 1\le qk+a\le X,\qquad pk+b_a+r\ge1,\qquad
 r\le s_a/q+\delta_X(qk+a)<r+1.
\end{equation}
Each set is an integer interval, possibly empty, because all its
restrictions are linear inequalities in $k$. For $c=0$ the last
restriction is either always true or always false.

The two affine arguments on this interval are
\begin{equation}\label{eq:forms}
 L_{1,a,r}(k)=qk+a,\qquad L_{2,a,r}(k)=pk+b_a+r.
\end{equation}
Their determinant is
\begin{equation}\label{eq:det}
 \Delta_{a,r}=q(b_a+r)-pa=qr-s_a.
\end{equation}
If it is zero, $qr=s_a$ forces $r=s_a=0$; coprimality then forces
$a=0$. Conversely, $(a,r)=(0,0)$ is proportional. This classification,
not merely subtraction of a proposed main term, defines the remainder:
\begin{equation}\label{eq:E}
 \mathcal E_{p,q,c}(X)=
 \sum_{\substack{0\le a<q,\ r\in\Z\\(a,r)\ne(0,0)}}
 \ \sum_{k\in I_{a,r}(X,c)}
 \lambda(qk+a)\lambda(pk+b_a+r).
\end{equation}
Only a fixed finite set of $r$ can occur. Indeed, for $n\le X$,
\[
 |\delta_X n|\le\frac pq |c|e^{|c|},
\]
using $|e^t-1|\le |t|e^{|t|}$. Thus the number of blocks depends on
$p,q,c$ but not on $X$.

Let
\begin{equation}\label{eq:N}
 N_{p,q,c}(X)=\#\{k\ge1:qk\le X,\ \lfloor pe^{c/X}k\rfloor=pk\}.
\end{equation}
\begin{proposition}[Exact structural decomposition: {\rm[LIT/ADAPT]}]
\label{prop:split}
With the preceding definitions,
\begin{equation}\label{eq:split}
 \boxed{\ C_X(u_X)=\lambda(pq)N_{p,q,c}(X)+\mathcal E_{p,q,c}(X).\ }
\end{equation}
Every summand defining $\mathcal E$ belongs to a nonproportional affine
block. The decomposition specializes to the residue-class decomposition
and the microscopic block partition of \cite[Sections 3 and 5]{resonance}.
\end{proposition}
\begin{proof}
Write $n=qk+a$. Since $pk+b_a$ is an integer,
\[
 \lfloor e^{u_X}(qk+a)\rfloor
 =pk+b_a+\lfloor s_a/q+\delta_X(qk+a)\rfloor.
\]
The intervals in \eqref{eq:I} partition precisely the allowed integers
in \eqref{eq:CX}. On the unique proportional block, $n=qk$ and the
second argument is $pk$, so its summand is
$\lambda(qk)\lambda(pk)=\lambda(pq)$. Its size is exactly \eqref{eq:N}.
The determinant calculation \eqref{eq:det} covers all other blocks.
\end{proof}

The task is to prove $\mathcal E_{p,q,c}(X)=o(X)$, rather than to
assume it through an affine Chowla hypothesis. The next sections develop
the required estimate from lower-level inputs in \cite{final}.

\section{Identify the exact analytic object in every remainder block}
\label{sec:conversion}

\begin{lemma}[Conversion to a progression interval: {\rm[DED]}]
\label{lem:conversion}
For a nonproportional block set
\[
 \ell=pq,\quad d_0=\min(pa,q(b_a+r)),\quad
 h=|q(b_a+r)-pa|>0.
\]
Then, for every $k$ in that block,
\begin{equation}\label{eq:conversion}
 \lambda(qk+a)\lambda(pk+b_a+r)
 =\lambda(pq)\lambda(pqk+d_0)\lambda(pqk+d_0+h).
\end{equation}
Consequently its sum is $\lambda(pq)$ times the sum of
$\lambda(m)\lambda(m+h)$ over the integer interval
$m\in pqI_{a,r}(X,c)+d_0$ restricted to $m\equiv d_0\pmod{pq}$.
The two arguments are positive there, even when $d_0<0$.
\end{lemma}
\begin{proof}
The two scaled arguments are $p(qk+a)=pqk+pa$ and
$q(pk+b_a+r)=pqk+q(b_a+r)$. They differ by $h$ and their smaller
one is $pqk+d_0$. Complete multiplicativity gives their Liouville
product as $\lambda(pq)$ times the product in \eqref{eq:forms}.
Multiplying again by $\lambda(pq)$ proves the identity. Positivity
follows from \eqref{eq:I}, since the scaled arguments are positive
multiples of the original arguments. The map $k\mapsto pqk+d_0$
is increasing and has image one residue class in an interval.
\end{proof}

We therefore need a bound for
\begin{equation}\label{eq:F}
 F_{h,\ell,b}(Y)=\frac1Y\sum_{1\le m\le Y}
 \1_{m\equiv b\pmod\ell}\lambda(m)\lambda(m+h)
 \quad(Y\ge1),
\end{equation}
with $h,\ell,b$ fixed. No coprimality condition on $b,\ell$ is available
or needed. We will derive the bound, rather than invoke the progression
conclusion or Theorem 1.1 of \cite{final}.

\section{Dilate the progression atoms and center the prime factors}
\label{sec:center}

Fix one triple $h,\ell,b$. For this section and the next two, use $Y$
as the long cutoff, and abbreviate $F=F_{h,\ell,b}$. Choose absolute
constants $A\ge1000$ and, later, $W\ge10$. Put
\begin{equation}\label{eq:parameters}
 \begin{gathered}
 L=(\log Y)^{1/A},\quad \delta=1/200,\quad
 J=\left\lfloor\frac{\delta\log L}{6W}\right\rfloor,\quad
 \eta=e^{-J},\quad K=e^{4J},\\
 U_i=L^{1-\delta}e^{6Wi}\ (0\le i\le J),\qquad
 P_i=\{s\text{ prime}:s\equiv1\pmod5,\ U_{i-1}\le\log s<U_i,
                         \ s\nmid h\ell\},\\
 Q=\{s\text{ prime}:s\not\equiv1\pmod5,\ s\le e^L,\ s\nmid h\ell\}.
 \end{gathered}
\end{equation}
The letter $s$ denotes a prime here; the rational parameters $p,q$
of Section~\ref{sec:remainder} remain fixed and separate. Set
\[
 V_i=\sum_{s\in P_i}s^{-1},\quad V_* =\prod_{i=1}^J V_i,\quad
 \mathcal D=\{s_1\cdots s_J:s_i\in P_i\}.
\]
The prime number theorem in the fixed modulus $5$, applied by partial
summation, gives $V_i=3W/2+o(1)$ uniformly in $i$. In particular,
\begin{equation}\label{eq:masses}
 W\le V_i\le2W,\qquad d\le e^{2L}\quad(d\in\mathcal D)
\end{equation}
for sufficiently large $Y$, with $h,\ell,W$ fixed. Let $\mathcal Q$
be the squarefree products of primes in $Q$, including $1$, and define
\[
 v(z)=4^{\omega(z)},\qquad
 S=\sum_{z\in\mathcal Q}\frac{v(z)}z=\prod_{s\in Q}(1+4/s),\qquad
 w(m)=\sum_{\substack{z\in\mathcal Q\\z\mid m}}v(z)=5^{\omega_Q(m)}.
\]
For integers $j\ge0$ with $j\eta\le100L$, put
\begin{equation}\label{eq:bins}
 \mathcal S_j=\{(d,z):d\in\mathcal D,\ z\in\mathcal Q,
  \ \omega(z)\le100\log L,\ j\eta\le\log(dz)<(j+1)\eta\},
 \qquad T_j=Ye^{j\eta}.
\end{equation}
There are $O(L/\eta)$ bins. Their retained harmonic mass is
\begin{equation}\label{eq:S0}
 S_0=\sum_j\sum_{(d,z)\in\mathcal S_j}\frac{v(z)}{dz},\qquad
 \frac12SV_*\le S_0\le SV_*.
\end{equation}
For the lower bound, under the law $v(z)/(zS)$, the inclusion probability
of a prime $s\in Q$ is $4/(s+4)$. Partial summation gives
$\E\log z\le4L$ and $\E\omega(z)\le4\log L$ for large $Y$.
Markov's inequality discards at most $4/98+4/100<1/2$ of that law by
requiring $\log z\le98L$ and $\omega(z)\le100\log L$. Equation
\eqref{eq:masses} puts all remaining pairs in the bins. The reciprocal
mass of the $d$ choices is $V_*$, proving \eqref{eq:S0}.

Define the centered correlation
\begin{equation}\label{eq:Cj}
 \begin{split}
 C_j=\frac1{T_j}\sum_{1\le m\le T_j}\sum_{(d,z)\in\mathcal S_j}
 &v(z)\1_{z\mid m}\1_{m\equiv bzd\pmod\ell}
 \lambda(m)\lambda(m+hzd)\\[-2pt]
 &\hspace{5mm}\cdot\prod_{s\mid d}(\1_{s\mid m}-1/s).
 \end{split}
\end{equation}
\begin{lemma}[The error of centering: {\rm[ADAPT]}]
\label{lem:center}
\begin{equation}\label{eq:center}
 \left|\sum_j C_j-S_0F(Y)\right|
 \ll_{h,\ell}SV_*\left(\eta+Y^{-1}
                           +2^J\eta^{-1}L^{-1/20}\right).
\end{equation}
\end{lemma}
\begin{proof}
Expand the product in \eqref{eq:Cj}. The term retaining all divisibility
indicators has $m=zdv$. Complete multiplicativity and invertibility of
$zd$ modulo $\ell$ turn this term into
\[
 \frac{v(z)}{zd}F(T_j/(zd)).
\]
Here $Ye^{-\eta}<T_j/(zd)\le Y$. Boundedness of the summands of $F$
therefore gives $F(T_j/(zd))=F(Y)+O(\eta+Y^{-1})$.
Summing yields the main term $S_0F(Y)$ and its first two errors.

For a partial term let $I\ne\varnothing$ be the subset of prime bands
whose constant terms were chosen. Write $d=d't$, where
$t=\prod_{i\in I}s_i$. Substituting $m=zd'v$ and extracting the
factor $(-1)^{|I|}v(z)/(zd')$ leaves
\begin{equation}\label{eq:roughmean}
 \frac1U\sum_{1\le v\le U}\sum_t\frac1t
   \1_{v\equiv bt\pmod\ell}\lambda(v)\lambda(v+ht),
 \qquad U=T_j/(zd').
\end{equation}
For fixed $z,d',j,I$, the possible $t$ occupy a multiplicative interval
of width $e^{\eta}<2$. Split it into at most two dyadic intervals and
the $\ell$ residue classes of $t$. On each part the progression
restriction on $v$ is fixed. Every prime factor of $t$ is at least
$\exp(L^{0.995})$, and the dyadic scale $D$ satisfies
$\tfrac12\exp(L^{0.995})\le D\le\exp(2L)$. A nonempty bin also gives
$U>Ye^{-\eta}t$, hence $U\ge\exp(L^A/2)$.

The Fourier and short-sum inputs are precisely Lemma 4.2 and Theorem
4.3 of \cite{final}, the latter quoting the short exponential-sum
theorem of Matom\"aki--Radziwi\l\l--Tao. They give, for
$B(\theta)=\sum_t\ee(ht\theta)/t$ on each such dyadic part,
\begin{equation}\label{eq:fourier-input}
 \|B\|_\infty\ll L^{-0.99},\qquad
 \int_\T|B|^4\ll D^{-1}L^{-3.96},\qquad
 \sup_\theta\sum_{v\le U}|F_v(\theta)|\ll UD L^{-0.8},
\end{equation}
where
\[
 F_v(\theta)=\sum_{a=1}^D\lambda(v+a)
             \1_{v+a\equiv b_0\pmod\ell}\ee(a\theta).
\]
In the third estimate the supremum is outside the average over $v$.
The progression restriction follows by its additive-character
expansion, whose coefficients have total absolute value $1$.
We import \eqref{eq:fourier-input}; we do not assert a proof here of
its sieve bounds or of the short exponential-sum theorem.

For completeness, its use in \eqref{eq:roughmean} is as follows.
Average over translations $a=1,\ldots,D$. The boundary error is
$O(D/U)$. Orthogonality expresses the remaining average as
\[
 \frac1{UD}\sum_{v\le U}\int_\T
 B(\theta)F_v(\theta)G_v(\theta)\,d\theta,
 \quad G_v(\theta)=\sum_{a=1}^{(2h+1)D}\lambda(v+a)\ee(-a\theta).
\]
On $|B|\le L^{-1.05}$, Parseval and Cauchy--Schwarz give a bound
$O_h(L^{-1.05})$. On the complement, the fourth moment gives measure
$O(D^{-1}L^{0.24})$. Using $|G_v|\le(2h+1)D$ and the last estimate in
\eqref{eq:fourier-input}, its contribution is
$O_h(L^{0.24-0.99-0.8})=O_h(L^{-1.55})$. Thus \eqref{eq:roughmean}
is $O_{h,\ell}(L^{-21/20})$.

For fixed $I$ and $j$, the total extraction weight is at most
$S\prod_{i\notin I}V_i\le SV_*$. Summing over fewer than $2^J$ sets
$I$ and $O(L/\eta)$ bins gives the last error in \eqref{eq:center}.
\end{proof}

\section{The graph input and the operator argument it supplies}
\label{sec:graph}

This section records exactly which additional parts of \cite{final}
are imported. They concern centered edges and finite residue tests;
none states ordinary two-point Chowla as a hypothesis.

First put
\[
 \rho_j(m,d)=\frac1{w(m)}
   \sum_{\substack{z\in\mathcal Q,\ z\mid m\\
             j\eta\le\log(zd)<(j+1)\eta}}v(z).
\]
This sum includes all padding divisors in its bin, without the bound
on $\omega(z)$ from \eqref{eq:bins}. Thus $\sum_j\rho_j\le1$.

Here is the deleted set from Definition 5.1 of \cite{final}, expressed
in the present notation. A step is $D_i=\varepsilon_i h z_i d_i$,
$\varepsilon_i\in\{-1,1\}$, with $(d_i,z_i)\in\mathcal S_j$.
A path is positive when $z_id_i$ divides the departure site of every
step. A forward prohibited word has length between $3$ and
$\lfloor L^{1/10}\rfloor$, unequal consecutive $d_i$, and consecutive
occurrences of each tuple prime. In addition, a prime $s\mid d_1$
absent from $d_m$ divides $\sum_{i=a}^mD_i$ for some $1<a<m$.
It is minimal if no shorter contiguous subword is forward prohibited
in either orientation. Let $\mathcal B_j$ be the starting sites of
positive paths with such minimal words. Put
\begin{equation}\label{eq:cutsdef}
 \begin{split}
 \tau_j(m,d)&=\1_{\omega_Q(m)\le400\log L}
       \1_{\rho_j(m,d)\le K/L}\1_{m\notin\mathcal B_j},\\
 \chi_j(m,d)&=\tau_j(m,d)\1_{\omega_P(m)\le6WJ},
       \qquad P=\bigcup_iP_i.
 \end{split}
\end{equation}
Let $C_j^\circ$ be \eqref{eq:Cj} with the extra factor
$\chi_j(m,d)\chi_j(m+hzd,d)$.

\paragraph{External deletion input {\rm[LIT]}.}
Lemma 5.3 of \cite{final} verifies the rare-event hypotheses for
$\mathcal B_j$; Lemma 4.6, using the finite integer-residue comparison
of Section 3 there, then states
\begin{equation}\label{eq:deletion}
 \frac1{SV_*}\sum_j|C_j-C_j^\circ|
 \ll_{h,\ell}2^J(K^{-1}+L^{-100}+e^{-2WJ})+e^{-L^{0.9}}.
\end{equation}
This is an integer-average estimate. Replacing the integer residues by
independent ones without the cited finite comparison would not justify it.

\paragraph{External moment and weight inputs {\rm[LIT]}.}
Set $M=\lceil e^{103L}\rceil$, $k=\lfloor L\rfloor$. On a block
$\{t+1,\ldots,t+M\}$ define the real symmetric matrix $A_d$ by the
following increasing edges, and their transposes:
\begin{equation}\label{eq:matrix}
 \begin{split}
 A_d(m,m+hzd)=
 &\frac{L v(z)\tau_j(m,d)\tau_j(m+hzd,d)}
          {\sqrt{w(m)w(m+hzd)}}\,
 \1_{z\mid m}\1_{m\equiv bzd\pmod\ell}\\[-2pt]
 &\hspace{10mm}\cdot\prod_{s\mid d}(\1_{s\mid m}-1/s),
 \qquad (d,z)\in\mathcal S_j.
 \end{split}
\end{equation}
All other entries vanish. Let $E$ be the coordinate projection to
$\omega_P(m)\le6WJ$, and on copies indexed by $d$ define
$H_{d,d'}=\1_{d\ne d'}A_{d'}$.
Theorem 5.4 of \cite{final}, proved there using Sections 5--7, supplies
an absolute $C_2$ such that
\begin{equation}\label{eq:moment}
 \operatorname*{avg}_{t\in I}\|H^k\|_{\rm HS}^2
 \le [K(C_2\sqrt W)^J]^{2k}
\end{equation}
whenever the integer interval $I$ of block origins has length at least
$\exp(L^A/2)$, with $A$ as required by that manuscript's finite-law
comparison. The threshold may depend on $h,\ell,W$.
Its Corollary 3.3 also supplies, uniformly in every offset $a$,
\begin{equation}\label{eq:weight}
 \operatorname*{avg}_{t\in I}
  w(t+a)\1_{\omega_Q(t+a)\le400\log L}\le2S.
\end{equation}
These two estimates are used for integer block origins; no independence
of Liouville values or of translated sites is assumed.

\begin{lemma}[Transfer for noncommuting edge operators: {\rm[ADAPT]}]
\label{lem:transfer}
Let $B_i$ be finitely many self-adjoint operators and $E$ an orthogonal
projection. Put $\mathscr H_{i,j}=\1_{i\ne j}B_j$. If $\|B_i\|\le a$
and $\sum_i\|B_iv\|^2\le b^2\|v\|^2$ for $v=Ev$, then
\begin{equation}\label{eq:transfer}
 \left\|E\Big(\sum_iB_i\Big)E\right\|
 \le3\max(a,b,\rho(\mathscr H)).
\end{equation}
\end{lemma}
\begin{proof}
This is Lemma 8.1 of \cite{final}; the short resolvent proof is included
to display the transfer being used. Put $m_0=\max(a,b,\rho(\mathscr H))$.
If $m_0=0$ all $B_i=0$. Otherwise take $t=(2m_0)^{-1}$. For real
$|u|\le t$, $I-u\mathscr H$ and all $I+uB_i$ are invertible. The
self-adjoint operator
\[
 Q(u)=I-\sum_i uB_i(I+uB_i)^{-1}
\]
is invertible too: a vector $v$ in its kernel would give
$z_i=(I+uB_i)^{-1}v$ with $u\mathscr Hz=z$, a contradiction.
Since $Q(0)=I$, it stays positive definite throughout $[-t,t]$.
The resolvent identity gives
\[
 Q(u)=I-u\sum_iB_i+u^2\sum_iB_i^2(I+uB_i)^{-1}.
\]
For $v=Ev$ the quadratic form of the last term lies between $0$ and
$2u^2b^2\|v\|^2$. Positivity at $u=t$ and $u=-t$ therefore yields
$|\langle v,\sum_iB_iv\rangle|\le(1+2t^2b^2)\|v\|^2/t
\le3m_0\|v\|^2$. Take the supremum over unit vectors in the range
of $E$. Distinct $B_i$ were not interchanged at any point.
\end{proof}

\begin{lemma}[The moment gives a retained operator bound: {\rm[ADAPT]}]
\label{lem:norm}
There is an absolute $C_3$ such that, outside a fraction at most
$e^{-2k}$ of admissible block origins,
\begin{equation}\label{eq:norm}
 \left\|E\Big(\sum_d A_d\Big)E\right\|
 \le K(C_3\sqrt W)^J.
\end{equation}
\end{lemma}
\begin{proof}
Put $a_d(m)=\prod_{s\mid d}|\1_{s\mid m}-1/s|$.
It is constant along an $A_d$ edge because that displacement is a
multiple of $d$. The definition of $\rho_j$ and the cutoff in
\eqref{eq:cutsdef} give the weighted row bound
\begin{equation}\label{eq:row}
 \sum_{m'}|A_d(m,m')|\frac{\sqrt{w(m')}}{\sqrt{w(m)}}\le2K a_d(m).
\end{equation}
For the incoming orientation, $z$ divides both endpoints or neither;
discarding the other endpoint restrictions gives the same bound.
The weighted Schur inequality applied on each invariant level set of
$a_d$ now gives
\[
 \|A_d\|\le2K,\qquad
 \|A_dv\|^2\le4K^2\sum_m a_d(m)^2|v_m|^2.
\]
At a site retained by $E$, arithmetic-geometric mean and
\eqref{eq:masses} give
\[
 \sum_d a_d(m)^2\le\sum_d a_d(m)
 \le\prod_i(\omega_{P_i}(m)+V_i)
 \le(8W)^J.
\]
Thus $\sum_d\|A_dv\|^2\le4K^2(8W)^J\|v\|^2$ for $v=Ev$.
Also $\rho(H)^{2k}\le\|H^k\|_{\rm HS}^2$. Markov's inequality
applied to \eqref{eq:moment} shows that
$\rho(H)\le eK(C_2\sqrt W)^J$ outside the stated exceptional
fraction. Lemma~\ref{lem:transfer} proves \eqref{eq:norm} with
$C_3=3\max(2,2\sqrt8,eC_2)$, independent of $W,h,\ell$.
\end{proof}

\section{Test the graph against the actual remainder atoms}
\label{sec:testing}

\begin{lemma}[Recover the finite prefix: {\rm[ADAPT]}]
\label{lem:testing}
For every bin,
\begin{equation}\label{eq:binbound}
 |C_j^\circ|\ll_{h,\ell}\frac{SK(C_3\sqrt W)^J}{L}+e^{-L}.
\end{equation}
\end{lemma}
\begin{proof}
Write $N=\lfloor T_j\rfloor$ and average over origins $0\le t<N$.
The length $N\ge\exp(L^A/2)$ is admissible in
\eqref{eq:moment}--\eqref{eq:weight}. On a block use the test vector
\[
 f_t(m)=\sqrt{w(m)}\lambda(m)
              \1_{\omega_Q(m)\le400\log L},\qquad g_t=Ef_t.
\]
Summing \eqref{eq:weight} over the $M$ offsets gives
$N^{-1}\sum_t\|g_t\|^2\le2MS$. On the good origins,
\eqref{eq:norm} bounds the average absolute quadratic form by
$2MSK(C_3\sqrt W)^J$. On every block the deterministic row bound
\eqref{eq:row} gives
\[
 \left|\left\langle g_t,\sum_d A_dg_t\right\rangle\right|
 \le2KM\,5^{400\log L}(8W)^J.
\]
For fixed $W$ its right side divided by $M$ is polynomial in $L$.
The bad fraction $e^{-2k}$ therefore contributes $O(e^{-L})$
after normalization by $2LM$. We have proved
\begin{equation}\label{eq:qform}
 \left|\frac1{2LMN}\sum_{t=0}^{N-1}
 \left\langle g_t,\sum_d A_dg_t\right\rangle\right|
 \ll\frac{SK(C_3\sqrt W)^J}{L}+e^{-L}.
\end{equation}

To justify that this controls the original prefix, write
$c_j(m;d,z)$ for its retained signed edge summand, so that
$C_j^\circ=T_j^{-1}\sum_{m\le N,d,z}c_j(m;d,z)$.
Testing an increasing entry of \eqref{eq:matrix} cancels both square-root
weights and gives $L c_j(m;d,z)$; the transpose gives the same value.
For $r=hzd$, the number of blocks containing both endpoints is
\[
 \mu_N(m,r)=
 [\min(N-1,m-1)-\max(0,m+r-M)+1]_+.
\]
Hence the expression inside the absolute value in \eqref{eq:qform}
is exactly
\[
 \frac1{MN}\sum_{m,d,z}\mu_N(m,hzd)c_j(m;d,z).
\]
All displacements are at most $R_L=he^{100L+1}$, and the one-orientation
row bound gives
\[
 \sum_{d,z}|c_j(m;d,z)|\le
 D_L:=\frac KL5^{400\log L}(8W)^J.
\]
For large $Y$, $N\gg M>R_L$. On $M\le m\le N$ one has
$\mu_N(m,r)=M-r$; the lower boundary and the extra sites $N<m<N+M$
contain at most $2M$ sites altogether. Thus replacing the averaged
quadratic form by $C_j^\circ$ has error at most
\[
 O\left(D_L\left(\frac{R_L}{M}+\frac MN+\frac1{T_j}\right)\right)
 =O(e^{-L}).
\]
The last equality uses $M=\lceil e^{103L}\rceil$, $T_j\ge Y=e^{L^A}$
and polynomial growth of $D_L$ for fixed $W$. This includes the real
cutoff rounding and terminal edges leaving the original prefix.
Combining with \eqref{eq:qform} proves \eqref{eq:binbound}.
\end{proof}

\begin{proposition}[The bound needed by every remainder block]
\label{prop:atoms}
Under the external intermediate inputs stated above, there is an
absolute $\kappa>0$ such that, for every fixed $h,\ell\ge1$ and
$b\in\Z$,
\begin{equation}\label{eq:atom}
 |F_{h,\ell,b}(Y)|\ll_{h,\ell}(\log Y)^{-\kappa}\qquad(Y\ge3).
\end{equation}
\end{proposition}
\begin{proof}
Use \eqref{eq:S0}, Lemma~\ref{lem:center}, the deletion estimate
\eqref{eq:deletion}, and sum Lemma~\ref{lem:testing} over
$O(L/\eta)$ bins. Since $V_*\ge W^J$, this gives
\begin{equation}\label{eq:finish}
 \begin{split}
 |F(Y)|\ll{}&\eta+2^J\eta^{-1}L^{-1/20}
 +2^J(K^{-1}+L^{-100}+e^{-2WJ})\\
 &+\frac K\eta(C_3/\sqrt W)^J+e^{-L^{0.9}}+Y^{-1}.
 \end{split}
\end{equation}
The accumulated $O((L/\eta)e^{-L})$ term is absorbed by
$e^{-L^{0.9}}$, as $SV_*\ge1$ and $L/\eta$ is polynomial in $L$.

Choose constants in the following order: first $A$ large enough for
the imported finite-law and moment statements; then their absolute
$C_2$ and the resulting $C_3$; finally $W\ge10$ so large that
$e^5C_3/\sqrt W\le e^{-1}$. All are fixed before $Y$ increases.
The spectral term in \eqref{eq:finish} is then at most $e^{-J}$.
The two elementary exponential errors obey
\[
 2^JK^{-1}=e^{-(4-\log2)J}\le e^{-J},\qquad
 2^Je^{-2WJ}\le e^{-J}.
\]
Since $J\le\delta\log L/(6W)$ and
$\delta(2+\log2)/(6W)<1/20$, one also has
$2^Je^JL^{-1/20}\le e^{-J}$; likewise $2^JL^{-100}\le e^{-J}$.
The last two errors are $O(e^{-J})$. Consequently
\[
 |F(Y)|\ll_{h,\ell}e^{-J}
 \le eL^{-\delta/(6W)}
 =e(\log Y)^{-\kappa},\qquad \kappa=\frac{\delta}{6WA}>0.
\]
This proves the estimate at every sufficiently large real cutoff.
The bounded remaining range is absorbed by enlarging the constant,
using $|F(Y)|\le1$. Residue representatives can be reduced modulo
$\ell$, so the constant need not depend on the representative $b$.
\end{proof}

\section{Close the original remainder, with all endpoint errors}
\label{sec:close}

\begin{theorem}[Quantitative closure of the resonance remainder]
\label{thm:closure}
Accept the external intermediate estimates of \cite{final} identified
in Sections~\ref{sec:center}--\ref{sec:graph}. For every fixed coprime
$p,q\ge1$ and fixed $c\in\R$, the actual block remainder
\eqref{eq:E} satisfies
\begin{equation}\label{eq:closure}
 \boxed{\ \mathcal E_{p,q,c}(X)
       =O_{p,q,c}\bigl(X(\log X)^{-\kappa}\bigr)=o(X).\ }
\end{equation}
The exponent $\kappa>0$ is absolute. No uniformity for $p,q,c$ growing
with $X$ is asserted.
\end{theorem}
\begin{proof}
There are finitely many nonempty block types, independently of $X$.
Fix one. Lemma~\ref{lem:conversion} identifies its sum as
$\lambda(pq)$ times a progression-interval sum with fixed $h,\ell,b$.
All its arguments $m=pqk+d_0$ are positive and at most
$pX+|d_0|$, hence at most $K_0X$ for a fixed $K_0$ and $X\ge1$.
The progression-interval sum is a difference of two prefixes
\[
 P_{h,\ell,b}(T)=\sum_{1\le m\le T}
 \1_{m\equiv b\pmod\ell}\lambda(m)\lambda(m+h),
 \qquad 0\le T\le K_0X,
\]
with the empty prefix defined to be $0$. If $T\le\sqrt X$, boundedness
gives $|P(T)|\le\sqrt X$. If $T>\sqrt X$, Proposition~\ref{prop:atoms}
and $\log T\ge\tfrac12\log X$ give
$|P(T)|\ll_{h,\ell,K_0}X(\log X)^{-\kappa}$.
As $\sqrt X=O(X(\log X)^{-\kappa})$, the same bound holds uniformly
for all endpoints in this fixed range. Two prefixes therefore bound
the block, without any assumption that its endpoints converge or that
its length is macroscopic. Summing the finite collection of blocks
proves \eqref{eq:closure}. Negative fixed intercepts cause no additional
issue, since the mapped arguments were already positive.
\end{proof}

This closes exactly the nonproportional terms left conditional in the
natural microscopic-profile proof of \cite[Section 5]{resonance}.
The proof did not use the source's final affine theorem. It did use
its intermediate graph-moment and deletion estimates, including their
finite-law justification. Removing those inputs would remove the
justification of \eqref{eq:closure}.

\section{Recover two-point Chowla from the negative remainders}
\label{sec:chowla}

This deduction makes the direction of the argument explicit:
we start with the resonance remainder, establish its cancellation,
and then extract ordinary fixed-shift cancellation.

For $j\ge1$ define, for real $T\ge0$,
\begin{equation}\label{eq:Sj}
 S_j(T)=\sum_{j<n\le T}\lambda(n)\lambda(n-j).
\end{equation}
\begin{lemma}[A negative microscopic remainder resolves terminal shift bands:
{\rm[DED]}]
\label{lem:bands}
For each fixed integer $h\ge1$,
\begin{equation}\label{eq:bands}
 \mathcal E_{1,1,-h}(X)
 =\sum_{j=1}^h
 \left[S_j\left(\frac{jX}{h}\right)
       -S_j\left(\frac{(j-1)X}{h}\right)\right]+O_h(1).
\end{equation}
For $h=1$ the stronger identity
$\mathcal E_{1,1,-1}(X)=S_1(X)$ holds exactly for every integer $X\ge1$.
\end{lemma}
\begin{proof}
Put $\alpha_X=1-e^{-h/X}>0$. For integer $n$,
\[
 \lfloor e^{-h/X}n\rfloor=n-\lceil\alpha_X n\rceil.
\]
There is no proportional block because $c=-h<0$, hence
$\mathcal E_{1,1,-h}(X)=C_X(-h/X)$. The correction is $-j$ precisely
when $(j-1)/\alpha_X<n\le j/\alpha_X$. Requiring the second argument
to be positive is precisely $n>j$, as encoded by $S_j$.
Since $\alpha_XX\to h$ from below, for sufficiently large $X$ the
possible $j$ are exactly among $1,\ldots,h$. We have the exact identity
\begin{equation}\label{eq:exactbands}
 \mathcal E_{1,1,-h}(X)=\sum_{j=1}^h
 \left[S_j\left(\min\left(X,\frac j{\alpha_X}\right)\right)
       -S_j\left(\frac{j-1}{\alpha_X}\right)\right]
\end{equation}
for such $X$ (the last lower endpoint is less than $X$).
Expansion at fixed $h$ gives
\[
 \frac1{\alpha_X}=\frac Xh+\frac12+O_h(X^{-1}).
\]
Each endpoint therefore differs by $O_h(1)$ from the corresponding
endpoint in \eqref{eq:bands}. Changing a bounded-summand prefix by
$O_h(1)$ changes it by $O_h(1)$, proving that formula.
For $h=1$, $0<\alpha_X n\le X(1-e^{-1/X})<1$, so the correction
is always $-1$. Only $n=1$ is omitted for a nonpositive second
argument, giving the exact identity claimed.
\end{proof}

\begin{theorem}[Ordinary two-point Chowla from remainder closure: {\rm[DED]}]
\label{thm:chowla}
If $\mathcal E_{1,1,-h}(X)=o(X)$ for every fixed positive integer $h$,
then
\begin{equation}\label{eq:chowla}
 \sum_{n\le X}\lambda(n)\lambda(n+h)=o(X)
 \qquad\text{for every fixed }h\ge1.
\end{equation}
In particular Theorem~\ref{thm:closure}, with its stated external
inputs, proves this assertion and gives the quantitative bound
$O_h(X(\log X)^{-\kappa})$. Conversely, \eqref{eq:chowla} for all
fixed $h$ implies $\mathcal E_{1,1,-h}(X)=o(X)$ for all such $h$.
\end{theorem}
\begin{proof}
We induct on $h$. The exact identity for $h=1$ in
Lemma~\ref{lem:bands} proves $S_1(X)=o(X)$.
Assume $S_j(T)=o(T)$ for $1\le j<h$. Every term with $j<h$ in
\eqref{eq:bands} is $o(X)$, so the assumed remainder cancellation gives
\begin{equation}\label{eq:tail}
 S_h(X)-S_h(aX)=o(X),\qquad a=(h-1)/h\in(0,1).
\end{equation}
This relation extends from integer to real $X$ by bounded endpoint
changes. Write $M_h=\limsup_{X\to\infty}|S_h(X)|/X\le1$.
The triangle inequality in \eqref{eq:tail} yields $M_h\le aM_h$;
since $a<1$, $M_h=0$. This proves the induction.
Equivalently, iterate \eqref{eq:tail} a fixed number of times, let
$X\to\infty$, and then let that number increase; the residual
prefix has relative length $a^r\to0$.

For the quantitative assertion, strengthen the induction to
$S_j(T)=O_j(T(\log T)^{-\kappa})$ for $T\ge3$.
Theorem~\ref{thm:closure} and \eqref{eq:bands} give
$S_h(X)-S_h(aX)=O_h(X(\log X)^{-\kappa})$ under that induction
hypothesis; the base case follows from the exact identity at $h=1$.
Iterate down to the first $r$ for which $a^rX\le\sqrt X$.
Every earlier scale has logarithm at least $\tfrac12\log X$, so the
sum of the errors is at most a constant times
\[
 \frac{X}{(\log X)^\kappa}\sum_{i=0}^{r-1}a^i
 \le \frac{X}{(1-a)(\log X)^\kappa}.
\]
The remaining prefix is at most $\sqrt X$, which is absorbed into
the same bound. This proves the quantitative induction.

Finally
$\sum_{n\le X}\lambda(n)\lambda(n+h)=S_h(X+h)$ for integer $X$,
so the proved prefix cancellation gives \eqref{eq:chowla}.
Conversely, \eqref{eq:chowla} gives $S_j(T)=o(T)$ for each fixed $j$;
inserting these estimates into the finite sum \eqref{eq:bands} proves
the converse.
\end{proof}

The negative microscopic remainders in this theorem contain the whole
fixed-shift two-point problem. Thus their $o(X)$ estimate is substantive
analytic work; it does not follow from the structural decomposition
alone. Here that work is supplied through the stated intermediate
inputs and the developed argument in Sections~\ref{sec:center}--\ref{sec:close}.

\section{The completed natural microscopic profile}
\label{sec:profile}

\begin{theorem}[The author's profile after closing its remainder]
\label{thm:profile}
With the same external inputs and fixed $p,q,c$ as in
Theorem~\ref{thm:closure},
\begin{equation}\label{eq:profile}
 \frac1X C_X\left(\log\frac pq+\frac cX\right)
 =\Psi_{p,q}(c)+O_{p,q,c}((\log X)^{-\kappa}),
\end{equation}
where the profile is the one in \cite[Theorem 5.1]{resonance}:
\[
 \Psi_{p,q}(c)=
 \begin{cases}
 0,&c<0,\\
 \lambda(pq)/q,&0\le c\le q/p,\\
 \lambda(pq)/(pc),&c>q/p.
 \end{cases}
\]
\end{theorem}
\begin{proof}
It remains only to count the proportional block, already classified
before introducing the remainder. For $c<0$, $pe^{c/X}k<pk$ for
every $k\ge1$, so $N_{p,q,c}(X)=0$. For $c=0$ it equals
$\lfloor X/q\rfloor$. For $c>0$ the exact condition is
$p(e^{c/X}-1)k<1$, and therefore
\[
 N_{p,q,c}(X)=\min\left(\left\lfloor\frac Xq\right\rfloor,
 \left\lceil\frac1{p(e^{c/X}-1)}\right\rceil-1\right).
\]
For fixed positive $c$,
\[
 \frac1{p(e^{c/X}-1)}=\frac X{pc}-\frac1{2p}+O_{p,c}(X^{-1}),
\]
whence $N_{p,q,c}(X)=X\min(1/q,1/(pc))+O_{p,q,c}(1)$.
This includes $pc=q$: equality at the limiting boundary does not
assert an exact finite plateau there. Combine this count with
\eqref{eq:split} and \eqref{eq:closure}, then divide by $X$.
The normalized $O(1)$ count error is absorbed into
$O((\log X)^{-\kappa})$.
\end{proof}

\paragraph{Logical and attribution boundary.}
The resonance block partition and the stated profile come from
\cite{resonance}. The Fourier, finite-law, deletion and moment estimates
identified above come from \cite{final} and its cited analytic inputs.
The remainder-to-progression conversion and the negative-band extraction
have been derived explicitly here. This establishes an application
of those inputs to the author's remainder; it provides no evidence
that the source manuscript historically used that decomposition.
It is also not a claim that the structural decomposition by itself
proves Chowla.

\paragraph{Route accounting.}
One route is developed: an \textbf{[ADAPT]} of the source's centered-graph
estimates to the author's remainder, with \textbf{[DED]} block conversion
and recovery of fixed-shift Chowla. No independent route outside the
source's analytic method is established or claimed.

\begin{thebibliography}{9}
\bibitem{resonance}
V. Casa\~na-Giner,
\emph{Rational Resonances and Microscopic Natural--Logarithmic Profiles
for the Liouville Function}, supplied author manuscript.
Source version: \path{tex_resonant_liouville_JNT_FINAL.tex}.
The exact rational partition is Proposition 3.1; the conditional
microscopic profile and its proof are Theorem 5.1 and Section 5.

\bibitem{final}
OpenAI, \emph{Ordinary two-point correlations of multiplicative functions},
supplied manuscript \texttt{final.pdf}, dated 24 September 2026.
External inputs used here: Corollary 3.3 (PDF page 12);
Lemma 4.2 and Theorem 4.3 (PDF pages 13--15);
Lemma 4.6 (PDF page 19); Definition 5.1 and Lemma 5.3
(PDF pages 21--23); Theorem 5.4 (PDF page 24, with its proof in
Sections 5--7). The operator transfer rederived here is Lemma 8.1
of that manuscript. Its Theorem 1.1 and Proposition 2.1 are not
used as inputs.
\end{thebibliography}
\end{document}
